\odiseachapter{Novena parte}{Hades}

Una prueba de que el cine puede captar, si se empeña, la complejidad de la literatura es la excelente versión que Nolan hace de la excursión al Hades. Las sombrías escenas del film parecen seguir al pie de la letra los versos de la \emph{Odisea}.

\begin{verse}
Y ya a la linde llegaba del hondicaudal Oceano.\\
Tienen los hombres cimerios allí su ciudad y poblado\\
en niebla y nubes envueltos; sin que jamás hacia abajo\\
a ellos el sol relumbrando los mire desde sus rayos\\
ni cuando pasa enhilando camino hacia el cielo estrellado\\
ni cuando vuelve otra vez a la tierra desde lo alto,\\
que solo fúnebre noche se tiende sobre esos cuitados.
\end{verse}

La mención a unas gentes ---los «hombres cimerios»--- que viven en el límite del mundo, en una noche continua, sugiere una atrevida especulación. ¿Y si esa Cimeria a la que se refiere Homero fuera la región polar, el extremo norte de Europa, donde la noche puede extenderse la mitad del año? ¿Y si el Poeta o la tradición bárdica que le precede hubieran tenido noticias del invierno en esa región desolada y remota y, no sin razón, la hubieran asociado a la antesala del averno?

Por otra parte, ¿qué antesala es esa? No parece que los navegantes hayan arribado al Hades propiamente dicho. No han cruzado la laguna Estigia ni pagado a Caronte, lo cual tiene sentido porque ninguno de ellos ha muerto todavía, así que, posiblemente, se están moviendo en una tierra de nadie donde pueden cohabitar con los muertos. Para invocarlos, hace falta, cómo no, sacrificar el ganado que han traído consigo:

\begin{verse}
Y Perimedes y Euríloco allí sujetaban las víctimas;\\
y yo, sacando en un vuelo de al muslo mi espada buída,\\
una hoya abrí como a un codo a lo largo y de abajo hasta arriba,\\
y las ofrendas por todos los muertos en ella vertía:\\
primero de lechimel, dulce vino al seguido, y, encima,\\
una tercera de agua; y espolvoreé blanca harina.\\
De esos los muertos con celo invoqué las cabezas baldías\\
con la promesa de hacer de una vaca horra y buena allá en Ítaca\\
un sacrificio en mi casa hartando de hornazos la pira,\\
y de apartar a Tiresias carnero que le inmolaría\\
solo para él, todo negro, de nuestro rebaño alegría.
\end{verse}

Se diría que la primera ofrenda es una especie de placebo que Odiseo acompaña de promesas de futuros sacrificios. El verso «De esos los muertos con celo invoqué las cabezas baldías» es espeluznante. Cuando ya los difuntos se han congregado a su alrededor, viene la matanza de verdad:

\begin{verse}
Cuando a las razas de muertos con mandas y con rogativas\\
tuve rezadas, cogiendo las reses les di degollina\\
sobre la zanja, y allí nubinegra la sangre corría:\\
desde el Erebo las sombras de muertos difuntos subían.\\
Novias, muchachos, ancianos con todo el sufrir de una vida,\\
niñas en flor que recién estrenaban en su alma las cuitas,\\
y muchos, muchos varones golpeados por asta broncina,\\
muertos en guerra y mojado de sangre su arnés todavía:\\
todas, cada una a su parte, en torno al hoyo bullían\\
en vocerío infernal; me agarró la pavura amarilla.
\end{verse}

Más espeluznante aún si cabe. La enumeración de los difuntos, la horrible democracia de la muerte, hace que se nos encoja el corazón como si estuviéramos viendo surgir de la tierra a los muchachos, a las niñas, a los ancianos y a las novias sin casar y a todos, todos esos guerreros muertos inútilmente. Por encima de todo, el griterío infernal que aterroriza incluso a Ulises, que tiene muy poco de pusilánime. Y Homero todavía no ha acabado:

\begin{verse}
Orden entonces les di ya a mis hombres para que enseguida\\
me desollaran las reses que allí degolladas yacían\\
a crudo bronce, y que me las quemaran, y que letanía\\
a Hades tremendo le hicieran y a Persefonea horripila;\\
mientras, sacando de nuevo de al muslo mi espada buída,\\
ahí me senté porque muerto ninguno su testa vacía\\
fuera a acercarla a esa sangre hasta haber de Tiresias noticia.
\end{verse}

Los aqueos desollan las reses a «crudo bronce» y Ulises, espada en mano, monta guardia para que ningún difunto «de testa vacía» beba antes que Tiresias. Volviendo al cine, todavía no he visto una película de zombis capaz de superar estas escenas escritas hace tres mil años. El primer difunto que aparece es Elpénor, el pobre muchacho que se ha matado unos días atrás en casa de Circe.

\begin{verse}
Y la primera en llegar fue la sombra de aquel nuestro Elpénor,\\
que de la tierra anchimunda seguía esperando sepelio,\\
y es que en la casa de Circe dejamos entonces su cuerpo\\
desnudo de honras y lloros, que urgíanos un otro duelo.\\
Y se llenó de tristeza mi alma y lloraba allí al verlo,\\
y le decía alzando la voz alada de verbos:

«Elpénor, ¿cómo has llegado a las sototerrañas tinieblas?\\
Más vía has hecho tú a pie que yo con mi nave renegra».
\end{verse}

El muchacho le suplica que le dé un entierro digno.

\begin{verse}
Y, cuando partas, atrás no me dejes sin honras ni duelos\\
abandonado, no vaya a venírsete la ira del cielo;\\
no, quémame, por favor, con mis armas, las todas que tengo,\\
y haz un montón para mí a la orilla del mar caniciento,\\
para que diga de un hombre infeliz a los bienvenideros.\\
Cúmpleme eso, y clava en el mismo montón ese remo\\
del que tiraba, estando yo vivo, con mis compañeros.
\end{verse}

Odiseo le asegura que así lo hará:

\begin{verse}
Dijo él así, y respondiéndole entonces de nuevo le digo:\\
«Eso por ti, pobre mío, lo haré; bien voy a cumplirlo».
\end{verse}

La escena nos reconcilia con el héroe milmañas, quizás por ser la primera vez que le vemos demostrando compasión, sufriendo genuinamente por el destino aciago de un pobre inocente. Pero a Odiseo aún le queda mucho sufrimiento por delante. La siguiente sombra en aparecer es su madre. Casi inmediatamente, Homero nos muestra que, compasión o no, Ulises es más duro de pelar que el bronce que esgrime:

\begin{verse}
Vino después de mi madre difunta el alma sombría,\\
la hija de Autólico cordialtanero, Anticlea, que viva\\
en casa yo la dejé cuando a Ilío marché la bendita.\\
Mi corazón se llenó de tristeza y lloré esa visita,\\
mas ni aun así permití que primera ---aunque bien me dolía---\\
la sangre fuera a tocar sin haber de Tiresias noticia.
\end{verse}

Hay que tener agallas para prohibirle a tu propia madre saciar la sed (de sangre) hasta que, por fin, aparece Tiresias y le dice al héroe:

\begin{verse}
«Sangre Laercia, Odiseo milmañas, hijo del cielo,\\
¿qué haces dejando la lumbre del sol, desgraciado, y viniendo\\
a este paraje sin goces por donde se ven solo muertos?\\
Ah, pero aparta del hoyo y aleja ya el filo del hierro\\
para que beba la sangre y te diga decir verdadero».
\end{verse}

Después de beber, Tiresias ofrece su profecía, confirmando los pesares que Odiseo va a sufrir por su arrogancia. Poseidón territodatremante le va a cobrar caro que cegara a su hijo. O para ser precisos, le va a cobrar caro que fuera tan idiota como para identificarse después de hacerlo.

\begin{verse}
Andas en busca del dulce retorno, preclaro Odiseo;\\
te lo pondrá difícil un dios; a hurtadas no creo\\
que del tremante tú puedas pasar; hay rencor en su pecho\\
airado porque a su hijo del todo dejástelo ciego.
\end{verse}

A pesar de lo cual, el ciego le da ciertas esperanzas:

\begin{verse}
Mas aun así aún podríais llegar, aunque penas sufriendo,\\
si tu deseo lo sabes frenar y el de tus compañeros\\
cuando te llegue la vez de arrimar el bien hecho velero\\
hasta la isla Trinacia, huyendo del mar violinegro,\\
y os encontréis campeando las vacas y gordos corderos\\
del Sol, que todo lo oye y todo lo ve desde el cielo.\\
Si sin tocar se las dejas y atiendes no más al regreso,\\
a Ítaca entonces podríais llegar, aunque penas sufriendo;
\end{verse}

En el film de Nolan, la granjera versión de Circe condena de antemano a los compañeros de Ulises. Su profecía no les da escape. Asegura que cuando lleguen a Trinacia, los aqueos sacrificarán las reses de Helios y eso les costará la vida. Pero Homero es más sutil y Tiresias deja abierta la posibilidad de salvación \emph{siempre que superen una prueba}, que, por otra parte, estamos seguros de que los hombres no superarán. Se trata de una vieja treta recurrente, el examen que, a priori, puede aprobarse pero siempre se suspende. Orfeo puede sacar a Eurídice del Hades, basta con que no mire atrás; Ulises y su gente podrían haber llegado sin más fatigas a Ítaca de no haber abierto la bolsa de los vientos; Perséfone solo tenía que abstenerse de comer la granada para no acabar en el Hades; Ícaro podía volar libremente siempre que no se acercara demasiado al sol. En todos los casos, la posibilidad de libre albedrío parece existir solo en teoría. La lectura del bienagudo director, condenando a los aqueos sin más juicio, no parece, por lo tanto, descabellada.

Tiresias sigue con su profecía:

\begin{verse}
pero si se las tocáis, desde aquí ya te anuncio hundimiento\\
para tu nave y tus hombres. Y aun cuando salieras tú ileso,\\
tarde, y no bien, llegarías, y todos tus hombres ya muertos,\\
y en nave de otros; y en casa otra peste vendría a tu encuentro:\\
machos soberbios que están devorándose vuestro sustento\\
de tu mujer divilinda en cortejo, y poniendo a ella precio.\\
Verdad es, sí, que al llegar vengarás las violencias de aquellos.
\end{verse}

Los versos, como muchos otros en la \emph{Odisea}, prefiguran ---¿justifican?--- la inevitable matanza. Después, Tiresias le pide a Ulises que peregrine hasta un lugar que vuelve a parecerse al fin del mundo y ofrezca un sacrificio para congraciarse con el enojado dios del mar.

\begin{verse}
Y cuando a los pretendientes al cabo por tus aposentos\\
los mates, ya con astucias, ya a bronce y al descubierto,\\
coge tu remo manobediente y camino ve haciendo\\
hasta que encuentres un pueblo de hombres que el mar nunca vieron,\\
ni la comida mezclaron con sal, y que nunca supieron\\
nada de naves rostribermejas y nada de remos\\
bienllevaderos, que mismo son alas para los veleros.\\
Señas muy claras te voy a decir, de fácil recuerdo:\\
cuando te encuentres, andando el camino, con otro andariego\\
que a ti te diga que al hombro gentil llevas palo de aviento,\\
hunde tu remo manobediente allí mismo en el suelo,\\
y a Poseidón soberano inmólale en aplacamiento\\
carnero, toro y verraco, y vuelve a tu casa de nuevo\\
para ofrecer una santa hecatombe a los no morideros,\\
los tantos dioses que pueblan la vasta campa del cielo,\\
uno tras otro hasta todos. Y a ti desde el mar ---estoy viendo---\\
muy fácil muerte te habrá de llegar, que el acabamiento\\
por apacible vejez abrumado vendrá; y con un pueblo\\
próspero siempre a tu lado. Y he dicho decir verdadero.
\end{verse}

¿Qué lugar es ese, alejado del mar, al que el héroe tiene que viajar para redimirse? Homero no da más explicaciones ---eso sí, volverá a repetir los versos, al pie de la letra, cuando Ulises se reencuentre con Penélope--- pero la tradición posterior le sacará jugo, como a casi todos los versos de la \emph{Odisea}. Lo interesante del pasaje es que prefigura un detalle que a menudo se pasa por alto. El viaje del héroe no termina en Ítaca, ya que está obligado a realizar un último viaje\footnote{¿Qué viaje es ese? Eugamón de Cirene (siglo VI a.C.) imagina en su \emph{Telegonía} que Ulises viaja a Tesprotia, en el Epiro, es decir, el noroeste continental de Grecia, frente a Corfú, más la franja meridional de Albania (el nombre Epiro significa literalmente «tierra firme»). La idea, recordemos, es expiar sus pecados, pero ya se sabe que Odiseo es incorregible y en lugar de purgar sus penas, se monta un segundo reinado, casándose con Calídice, reina de los tesprotos, con la que tiene un hijo, Polipetes. En cuanto a la profecía de Tiresias, «Y a ti desde el mar ---estoy viendo--- muy fácil muerte te habrá de llegar», el agente no es otro que Telégono, el hijo de Odiseo y Circe, que lo mata por error con una lanza rematada con aguijón de raya.}.

Dicha su profecía, Tiresias se retira y Ulises permite que su madre beba la sangre, lo que le devuelve ---quizás temporalmente--- el entendimiento. Anticlea le pregunta por sus andanzas y Odiseo, después de explicarse, le pregunta:

\begin{verse}
[...] Te pido\\
que hables tú ahora y que venga verdad verdadera a mi oído:\\
¿qué malhadar del morir sombriluengo hizo doma contigo?\\
¿Un largo mal, o la milflecheríos Artemis te vino\\
a visitar con sus dardos de dulce pasión de un suspiro?\\
Dime del padre, y del niño que yo abandoné, de mi hijo, y\\
si aquel mi reino aún con ellos está, o ya otro distinto\\
de entre los hombres lo tiene, y a mí ya me dan por perdido.\\
Dime también el querer y el pensar que mi esposa ha tenido,\\
si está ella aún con el hijo y cuida de todo lo mismo\\
o la llevó por mujer otro aqueo que más le convino.
\end{verse}

Anticlea le da noticias de Penélope, de Telémaco y de su anciano padre Laertes y le confiesa que la razón de su muerte fue la tristeza por el hijo ausente. Ulises intenta abrazarla en vano, un pasaje todavía más desgarrador que todos los desgarradores versos que ofrece el canto más triste de la \emph{Odisea} y puede que de toda la literatura, hasta que Primo Levi escribiera \emph{Si esto es un hombre}:

\begin{verse}
Eso me dijo, y yo en mis adentros ahí me moría\\
por agarrar de mi madre difunta aquel soplo de vida.\\
Tres veces a ella me fui ---a ella a mí el corazón me movía---, y\\
dejó, tal sombra o ensueño, tres veces mis manos vacías.\\
Más me punzaba el dolor cada vez, y la melancolía,\\
y, levantando de verbos alada la voz, le decía:

«Madre, ¿por qué no te paras ahí cuando quiero agarrarte\\
para poder envolverte en mis brazos, por más que en el Hades,\\
hasta saciarnos los dos en el gélido llanto? ¿O imagen\\
eres no más que Perséfone altiva ha querido mandarme\\
para que llore yo más todavía y el planto no acabe?».
\end{verse}

Anticlea le responde:

\begin{verse}
«Ay, hijo mío, que de hombres tú llevas la más triste estrella,\\
no te ha engañado la hija de Zeus, Persefonea,\\
que de mortales es esta la ley cuando muerte les llega;\\
a carne y huesos entonces no hay ya nervazón que los tenga,\\
pues del ardiente fuego la rinde la puja soberbia\\
en cuanto el soplo de vida abandona la blanca osamenta\\
y el alma, como si sueño, revolotea y se aleja.\\
Vuelve cuanto antes tú ahora a la luz, y porque un día puedas\\
esto a tu esposa contar, lo que estás aquí viendo recuerda».
\end{verse}

Permíteme, amable lectora, comprensivo lector, que en este punto me adelante al relato que Odiseo les está haciendo a los feacios. Después de hablar con su madre, aparece en el canto XI una larga lista de ilustres difuntas y más tarde el ínclito ---es un decir--- Agamenón, al que dedicaremos el próximo capítulo. Aquí, como buenos aficionados al cine, vamos a adelantar la película hasta la escena en la que aparece nada menos que Aquiles, junto con algunos otros héroes de la \emph{Ilíada}\footnote{Escena que también se salta Nolan. ¡Ay, las exigencias del guion!}. Ulises todavía está conversando con Agamenón cuando las sombras de los guerreros se acercan:

\begin{verse}
Fue esa la escena: plantados los dos y en odiosa conversa\\
---lágrimas mil y de luto---, en esas allí se presenta\\
el alma en pena de Aquiles Pelida, y al lado de ella\\
la de Patroclo, y también la de Antíloco prosapiabuena,\\
y la de Áyax, que en porte era él el mejor y en belleza\\
de entre los dánaos, sacando al Pelida sin par de esa cuenta.\\
Me conoció del Eácida el alma, pie raudo como era,\\
y entre sollozos palabras aladas me da que dijeran:

«Sangre Laercia, Odiseo milmañas, hijo del cielo,\\
¿a qué, gran terco, revuelve una hechura mayor aún tu pecho?\\
¿Cómo has osado bajar aquí al Hades, donde desatentos\\
viven los muertos, figuras no más de los hombres que fueron?».
\end{verse}

¡Cuánta razón tiene el Pelida! Pasearse por el infierno no es cosa baladí y hay que estar muy motivado para moverse entre «figuras no más de los hombres que fueron». El tema de las sombras se repite una y otra vez. Ni siquiera el más grande de todos los guerreros se escapa a ese terrible destino. Ulises se justifica y trata de animarle:


\begin{verse}
Eso me dijo, y respuesta le daba yo a él enseguida:\\
«Aquiles el de Peleo, el de más en los griegos valía,\\
necesidad de Tiresias me trajo, por si él me daría\\
una palabra que al fin me llevara a las breñas de Ítaca.\\
Que aún no he llegado ni cerca de Acaya, ni pie puse encima\\
de esa mi tierra, que siempre ando en mal, pero, Aquiles, tu dicha\\
no la vistió antes que tú hombre ninguno, ni habrá quien la vista.\\
Las mismas honras a ti que a los dioses te hicimos en vida,\\
todos los griegos, y ahora de muertos la soberanía\\
tienes aquí: tu morir, Aquileo, que en nada te aflija».
\end{verse}

La respuesta de Aquiles es devastadora:

\begin{verse}
Eso le dije, y respuesta me daba con esto en un vuelo:\\
De mi morir no me busques consuelos, preclaro Odiseo.\\
Preferiría doblarme en la tierra para otro y ser siervo\\
de uno sin casi ni hacienda y muy escaso sustento,\\
a ser el rey entre toda esta plebe acabada de muertos...
\end{verse}

Si necesitábamos una prueba definitiva de los horrores del Hades, aquí la tenemos. Aquiles preferiría ser el más humilde de los siervos bajo el luminoso sol que un rey de ultratumba. Y como el resto de los pobres difuntos, su única obsesión es que le den noticias de los vivos. Pregunta a Odiseo por su hijo y por su padre. Milmañas no sabe nada de este último, pero le refiere todas las aventuras ---es decir, las matanzas--- de su muchacho Neoptólemo, otro personaje adorable, muy amigo de destriparte mientras te saluda. Con esto, Aquiles se marcha contento:

\begin{verse}
Eso le dije, y el alma del hijo de Eaco, pie raudo,\\
por la pradera de asfódelos ya se perdía en dos trancos\\
alegre de ese hijo suyo que supo por mí tan preclaro.
\end{verse}

Siguen unos versos muy curiosos, que demuestran que los argivos son tan orgullosos que ni siquiera la muerte sirve para que se les pasen los enfados. Áyax todavía se la tiene jurada a Odiseo\footnote{Con cierta razón. La historia se relata con detalle en la \emph{Ilíada menor} de Lesques y suele recibir el nombre de \gr{ὅπλων κρίσις} (hóplōn krísis), el juicio de las armas ---en la expresión podemos ver la raíz de la palabra «hoplita», que designa al soldado griego, y de la moderna palabra «crisis»---. Muerto Aquiles, sus armas se ofrecen como premio. Áyax las reclama por haber rescatado el cadáver del héroe; Odiseo alega sus muchas contribuciones a la guerra y se queda con el premio. Áyax enloquece, degüella el ganado del ejército creyendo que eran los Atridas y, cuando recobra la razón, se mata con la espada que le había regalado Héctor.}:


\begin{verse}
Solo del Telamoníada el alma sola, de Áyax,\\
se mantenía distante, airada porque lo ganara\\
en el certamen que al pie de las naves movieron las armas\\
de Aquiles; premio a cobrar que su madre fijó veneranda.\\
Teucros y Palas Atena la justa aquella arbitraban.\\
Ay, ojalá nunca hubiera ganádole yo esa batalla;\\
que fue por ellas por lo que la tierra cubrió testa tanta:\\
Áyax, que en obra el primero, y también el primero en estampa\\
era de todos los dánaos, después del Pelida sin tacha.
\end{verse}

Es un destino dramático, el de este héroe, uno de los más simpáticos de la \emph{Ilíada}. Una verdadera torre ambulante se enfrenta con Héctor en una batalla cuerpo a cuerpo que acaba en empate técnico, aunque, a decir verdad, el héroe troyano estaba recibiendo una fenomenal paliza. Pero después de combatir, ambos enemigos se intercambian regalos, como nobles guerreros que son. Áyax le regala a Héctor un cinturón ---que luego usará Aquiles para arrastrarlo con su carro alrededor de las murallas de Ilión--- y el troyano le regala al gigante argivo su espada, que sirve para que este se suicide. Tragedia sobre tragedia, incluso para los más nobles.

Ulises intenta congraciarse con él y le dice:

\begin{verse}
«Áyax, retoño del buen Telamón, ¿no has pensado, ni muerto,\\
arrinconar ya la saña por esas armas, que hoy siento\\
armas malditas? Fue ruina que quiso de argivos el cielo,\\
pues tú moriste y tú eras su torre; por ti los aqueos,\\
tal por la santa cabeza de Aquiles el de Peleo,\\
siempre de muerte penando nos vemos; culpable de eso\\
solo lo es Zeus que a la hueste de dánaos astilifieros\\
la aborrecía a mayores, y el lance te impuso postrero.\\
Ven, ven aquí, gran señor, y da oído, que hablarte yo quiero,\\
a mis palabras; y mata esa rabia y arrojo altanero».
\end{verse}

Pero si a Áyax no le gana nadie en lo noble, tampoco en lo rencoroso:

\begin{verse}
Le dije. No contestó, y de las almas de muertos difuntos\\
hacia el Erebo buscó compañía. Sé bien que a él, incluso\\
enrabietado, debí hablarle más, y que hablara él lo suyo…
\end{verse}

Pero estamos en el Hades, lo que quiere decir que lo que no se arregló en vida tampoco se puede arreglar en la muerte. Las sombras carecen de voluntad y nada puede ya cambiarse. El canto concluye con una antología de notables difuntos que incluye a Minos, Orión, Tántalo y, cómo no, el célebre Sísifo:

\begin{verse}
Y pude a Sísifo ver, que sufría muy grave martirio\\
teniendo a peso en sus manos aquel pedrusco grandísimo.\\
A él apoyado con manos y pies, intentaba subirlo\\
por la pendiente de una montaña; y cuando iba ya mismo\\
a tramontar, por su peso el pedrusco hacia atrás de corrido\\
a la llanura rodaba otra vez, el maldesabrido.\\
Y se estiraba otra vez a empujarlo, y bajábanle ríos\\
de resudor por la piel, su cabeza de polvo en un nimbo.
\end{verse}

Concluye el desfile nada menos que Hércules; es decir, su imagen (otro \gr{εἴδωλον}, o proyección astral, como el que Hera quizás manda a Troya en lugar de Helena).

\begin{verse}
Luego de él, mis ojos los puse en Heracles violento,\\
bueno, en su imagen; pues que él acompaña a los dioses eternos,\\
junto a Hebe tobillolindo, al goce de los festejos,\\
la joven hija de Hera sandaliadeoro y de Zeus.\\
De su fantasma, tal grullas, salía un estruendo de muertos\\
tolondros y en desbandada; y él, en negrura semejo\\
a la honda noche, con su arco desnudo y la flecha en el nervio,\\
acá oteando y allá, pronto ya a disparar, daba miedo.
\end{verse}

Sombras que repiten los ecos de los vivos. Áyax eternamente rencoroso, Heracles eternamente violento. Los dioses, de hecho, se han molestado en fabricar una sombra de Heracles solo para que haga guardia en el Hades, mientras el original mora con los olímpicos. No es de extrañar que, muchos años más tarde, el anciano Ulises se pregunte si ese lugar horrible que creyó visitar existe realmente\footnote{\emph{Abandonando Ítaca}, \url{https://eolasediciones.es/libro/abandonando-itaca-leaving-ithaca/}.}:

\begin{verse}
Tú, que ya caminaste por los campos baldíos\\
y hablaste con las sombras que habitan sus tinieblas,\\
acuérdate de Aquiles, que gustosamente habría cambiado\\
su espada fantasmal por una humilde azada\\
de cualquier campesino. No te hagas ilusiones\\
sobre lo que te espera. Sin embargo, tu\\
corazón se enfurece ante la perspectiva.\\
Ser tan solo un espectro sin cerebro, sediento de sangre,\\
desprovisto de fuerza y de pasión, ¡vacío de deseo!\\
¡Ah, si la muerte fuese simplemente una luz que se apaga,\\
como algunos filósofos defienden con obstinación!

A veces uno se pregunta\\
si tu viaje al Hades fue real o tan solo\\
una pesadilla invocada por los filtros de Circe.\\
Tal vez, con suerte, no haya vida después de morir.\\
Sí, morir y no ser, disolverse en la nada,\\
olvidar los lamentos, desechar para siempre\\
los fantasmas de ese niño inocente,\\
de esos pobres muchachos.\\
Aferrarte, hasta el último instante,\\
a lo que fue bueno en tu vida. El primer beso\\
perdido entre naranjos, su cuerpo menudo ardiendo contra el tuyo.\\
La imagen de muchachas bailando desnudas en la playa,\\
y los inmensos ojos de tu hijo,\\
ofreciendo la última redención.
\end{verse}

En nuestro propio tiempo, obsesionados como estamos con la vida eterna que se nos sigue negando igual que a los griegos y con la eterna juventud que ni el bótox ni la cirugía plástica pueden comprar ni siquiera a los más pudientes, la idea de que la muerte pueda ser redención se nos hace extraña y casi repugnante. La culpa es de los dos mil años que llevamos oyendo hablar del Paraíso y la resurrección. No es que nadie se tome muy en serio esas leyendas hoy en día, pero sí nos han hecho olvidar que para los helenos, la vida después de la muerte no era vida, sino una parodia grotesca, un horrible existir sin existir. De ahí que Demócrito declare que la muerte son buenas noticias, de ahí que Epicuro nos diga que no hay que temerla ni preocuparse por ella ya que nos liberará, para siempre, del Hades.
